\section{Problem Formulation}
\label{sec:problem}
We retain the prelaunch allocation problem introduced in
WORKS~\cite{camp2026works}. A workflow comprises named processes, each of
which instantiates tasks for different samples or input partitions.
Before task $t$ starts, the allocator issues an initial memory request
$\hat y_t$. Its static feature vector $\mathbf{s}_t$ contains workflow and
process metadata, software and system context, resource settings, and
staged input size. Execution subsequently reveals peak resident set size
(RSS), $y_t$, and duration $d_t$. If the task is audited, it also supplies a
consumption measurement $c_t$. Neither $y_t$ nor $c_t$ is available for that
task's initial allocation.

Let $\mathcal{P}(t)$ contain tasks whose outcomes have been recorded before
the allocation decision for $t$. Peak-memory history $\mathbf{h}_{m,t}$
summarizes compatible observations from this set; consumption history
$\mathbf{h}_{c,t}$ uses its audited subset. The prelaunch decisions are
\begin{equation}
\hat y_t=f(\mathbf{s}_t,\mathbf{h}_{m,t},\mathbf{h}_{c,t}),
\qquad q_t\in\{0,1\},
\label{eq:allocation_problem}
\end{equation}
where $q_t=1$ selects auditing. A budget $B\in(0,1]$ limits the audited
fraction over the specified evaluation horizon. Auditing task $t$ cannot
improve the request already issued to it; its measurements can inform
only later decisions. In live execution, availability follows completion
and ingestion times, rather than launch order. Offline evaluation must
state the history construction and outcome-revelation protocol used to
approximate this constraint.

The allocation objective is to limit underallocation while avoiding
unnecessary reservations. For $N$ evaluated tasks, initial-request
coverage is
\begin{equation}
\mathrm{Cov}=\frac{1}{N}\sum_{t=1}^{N}\mathbf{1}[\hat y_t\geq y_t],
\label{eq:coverage}
\end{equation}
where $\mathbf{1}[\cdot]$ is the indicator function. This is the
\emph{first-attempt coverage} measure used in WORKS: retry increases do not
enter $\hat y_t$. A request below a recorded peak is an allocation
diagnostic, not proof of an observed OOM failure. Live studies therefore
report task completion, failures, and retries separately.

Coverage is considered alongside aggregate requests $\sum_t\hat y_t$ and
unused memory-time,
\begin{equation}
W=\sum_{t\in\mathcal{D}}\max(\hat y_t-y_t,0)\,d_t,
\label{eq:unused_memory_time}
\end{equation}
where $\mathcal{D}$ identifies tasks with valid duration measurements.
Using memory in GiB and duration in hours gives $W$ in GiB-hours. This
peak-based measure quantifies reservation above observed peak demand;
it is not an integral of instantaneous unused RSS. Aggregate requests
likewise do not represent simultaneous cluster usage. Missing durations
and retry accounting must be stated rather than treated as measured zero
cost. When estimating net savings, auditing overhead must be charged in
compatible units and distinguished from prediction or retraining cost.

For the live study, both new-input branches update their models, history,
and calibration from completed outcomes. They differ in whether a further
pass over the original inputs precedes the new workload. This tests the
value of additional familiar-input feedback during subsequent adaptation;
it does not isolate the effect of updating relative to a frozen allocator.
Both branches exclude the current task's outcome from its prelaunch decision.
